% ============================================================
%  Georg Brandes, Samlede Skrifter. Trettende Bind (Supplementbind).
%  Kjøbenhavn: Gyldendalske Boghandels Forlag (F. Hegel & Søn).
%  Græbes Bogtrykkeri.  1903.
%  Five pieces on Rasmus Nielsen and his school, 1867.
%  TRANSCRIPTION (Danish), printed pp. 85--109.
%
%  One of three transcription directories cut from this volume:
%      striden/               pp. 1--42    Indledning; Striden om Tro og Viden
%      nielsen-anmeldelser/   pp. 85--109  five pieces on Nielsen, 1867
%      brochner-bidrag/       pp. 115--123 review of Brøchner's Bidrag (1869)
%  NOT transcribed: pp. 43--84 (Dualismen i vor nyeste Filosofi), 110--114,
%  124 ff.  See the note on pp. 43--84 below.
%
%  ---- DATE -------------------------------------------------------------
%  The title page (PDF 7) reads 1903.  The Indledning (p. 3) is signed
%  "April 1903." / "G. B."  (catalog.yaml's 1899 is the series date, not
%  this volume's.)  This is Brandes's 1903 revision of pieces first printed
%  1865--69, in the orthography of 1903 (Filosofi, Teologi, eksakt, hv-,
%  single vowels), transcribed as printed.  Nothing has been harmonised
%  with the first printings or with the 1866 Dualismen.
%
%  ---- SOURCE -----------------------------------------------------------
%  bibliotek/Brandes, Georg/1903-samlede-skrifter-13.pdf
%      Google Books (University of Chicago copy), 552 PDF pp., 600 dpi
%      bitonal JBIG2, embedded text layer (witness 1).
%      sha256 9579f1fd8672884bcc7e149b43f5c8c640dcaaefa1aceca01d5f89cbed016704
%  Witness 2 (never copy-text): Project Runeberg's OCR text of the same
%      edition (gbsamskr/13), bibliotek/Brandes, Georg/gbsamskr-13.txt,
%      cut into printed pages by alignment against witness 1 (per-page word
%      agreement between the two witnesses, pp. 85--109: mean 0.974, min 0.951 over 25 pages).
%      It is a different OCR engine's reading; whether Runeberg photographed
%      a different physical copy is not known.
%
%  ---- PAGE MAP ---------------------------------------------------------
%      printed p. = PDF p. - 12, uniform over printed 1--136.
%  Verified 2026-09-22 AT THE IMAGE: folio read off the rendered page at
%  printed 9, 29, 38, 42, 44, 86, 97, 104, 109, 116, 123 (all match; pp. 2
%  and 17 print no folio and were placed by content).  Independently read
%  by H. Halvorson at 8, 9, 42, 43, 54, 85, 93, 123, 136.
%
%  ---- LAYOUT OF THE VOLUME ---------------------------------------------
%  Every piece begins on a fresh page with a sunk head: title in large
%  italic, the year of first publication in roman in parentheses, a short
%  rule; the opening page carries no folio and no running head.  Every piece
%  ends with a short centred rule (\piecerule).  Running heads carry a short
%  title of the piece (a third reading, recorded at each piece's head where
%  it differs); they are not transcribed.
%
%  ---- CONVENTIONS ------------------------------------------------------
%  Diplomatic.  Orthography, punctuation and capitalisation as printed.
%  Emphasis: ITALIC is the volume's only emphasis device -- no letterspacing
%    and no bold were found in pp. 1--123 -- and every italic span is set
%    \emph{}, whether it marks emphasis, a title or a foreign word; no
%    attempt is made to distinguish the functions.
%  Quotation marks: the volume quotes with SINGLE guillemets pointing
%    outward, ‹...› (U+2039/U+203A), for quotations, titles and terms alike;
%    double «...» occur only in a head (p. 102).  Transcribed as printed.
%  Dashes: spaced em-dash, set ---.  ɔ: as printed.
%  Footnotes: *), set with \fnmark{*)}{...} at the mark.
%  A word broken across a printed line is joined (Cau\-%); a hyphen printed
%    at a line end as part of the word is kept.
%  Printer's errors are transcribed AS PRINTED and logged in a % comment at
%    the site; nothing is silently corrected.
%  Punctuation is as the Chicago copy prints it.  The 600 dpi bitonal image
%    sometimes loses the tail of a faint comma; where it shows a point but
%    witness 1 (Google's OCR, made from its capture BEFORE compression to
%    bitonal JBIG2) reads a comma, the comma is transcribed and commented;
%    where the image and witness 1 both show a point, the point is
%    transcribed and logged as a printer's error, whatever the syntax wants.
%  Letters the bitonal scan may not resolve (ø/o, æ/œ, worn e/c): where the
%    printed form would not be a word, the word is read unless the glyph
%    positively shows otherwise (TRANSCRIPTION-PLAYBOOK §2 rule 3); each case
%    is commented.  In this fount the italic æ has a closed o-like bowl and no
%    separate a-stem, so it looks like œ (see p. 27 `duæ', p. 10 Fædrelandet).
%  A hyphen at a line end that could be either a compound hyphen or a word
%    break is resolved by the volume's own usage MID-LINE: printed only solid
%    -> joined; only hyphenated -> kept; both or never -> the printed hyphen
%    is kept.  Each case is commented with its evidence.
%  A word divided across a PAGE turn is joined with `%' before the page
%    marker (e.g. Spørgs%  / \opage{122}maal); the divisional hyphen is not set.
%
%  ---- CONTENTS vs HEADS ------------------------------------------------
%  The Indhold (PDF 9) gives short titles; the heads over the pieces name
%  the book under review and the year.  Both readings are kept: the Indhold
%  reading is the first argument of \piecehead (table of contents, running
%  head), the head as printed is the second.  Recorded at each piece.
%
%  ---- READINGS OF THE HEADS ----------------------------------------------
%  Indhold (PDF 9)                       Head as printed over the piece
%  Om den gode Vilje som Magt i          R. Nielsen: Om den gode Vilje som Magt i
%    Videnskaben 85                        / Videnskaben (1867)
%  Nielsen og Grundtvig 93               Gb. (Rudolf Schmidt): / R. Nielsens
%                                          Filosofi og den Grundtvigske
%                                          Anskuelse. (1867).  [stop after year]
%  Et Svar 98                            Et Svar (1867)
%  Fædrelandets Hr. 99 102               «Fædrelandet»s Hr. 99 (Bjørnstjerne
%                                          Bjørnson) (1867) -- the volume's
%                                          only double guillemets
%  R. Nielsens Lære om Tro og Viden 106  Heegaard: Professor R. Nielsens Lære om
%                                          / Tro og Viden (1867)
%
%  ---- PRINTER'S ERRORS AND DOUBTFUL READINGS (each commented at its site) --
%  p. 87 stray comma `bærer sig, saaledes'.  p. 90 `Gennemgangsmedium,'
%  before capital `Idet'; `gør' in small type, slash faint (sense reading).
%  p. 91 ‹Skal Strauss ... never closed.  p. 92 wrong-fount o in `Tro'.
%  p. 95 closing guillemet turned: stimle‹.  p. 96 `føler', damaged ø.
%  p. 98 `Hr. Gb en' (no stop).  p. 99 `fil at se'; `at samme Art'.
%  p. 104 `Kultur›?›' (extra closing mark); `af den ‹danske Kultur›'.
%  p. 107 `ham. siden'.
%  Line-end hyphens: `religions-/filosofiske' (pp. 96, 102) kept -- the
%  volume prints both forms mid-line; `Non-/sens' (p. 103) joined -- solid
%  mid-line, p. 88; `Halv-/Grundtvigianerne' (p. 101) and
%  `Natur-/Oprindelighed' kept.
%  Guillemet balance ‹ minus › = +2, and every unit of it is one of the
%  logged defects above (p. 91 +1, p. 95 +2, p. 104 -1).
%
%  ---- pp. 43--84, NOT TRANSCRIBED --------------------------------------
%  pp. 43--84 reprint "Dualismen i vor nyeste Filosofi" (1866).  The 1866
%  first edition is already in this collection (texts/brandes/dualismen),
%  fully transcribed and collated against two scans.  The 1903 text is
%  Brandes's own revision, in modernised orthography and with altered
%  chapter titles (Indhold: I. Indledning 43, Den absolute Uensartethed 54,
%  Det Antirationelle 67, Dualismens Følger 71), so it is a VARIANT WITNESS
%  to the 1866 text and deserves a collation of its own, not a second
%  transcription.
%
%  ---- HOW THIS TEXT'S ACCURACY WAS ESTABLISHED ------------------------
%  Word-accuracy: every batch diffed against TWO machine witnesses before
%  splicing (ocrdiff.py; TRANSCRIPTION-PLAYBOOK §3 step 4) -- witness 1 the
%  embedded layer, witness 2 Runeberg's OCR -- and then checked with
%  consensus.py, which compares case-sensitively and with punctuation where
%  the two witnesses agree with each other (ocrdiff lower-cases and ignores
%  punctuation).  3 batches; witness 1: 6 candidates, witness 2: 17,
%  in both: 1, consensus: 3; 0 transcription errors surfaced, all were the
%  witnesses' faults; 0 unresolved.
%  Independent collation: a seeded random sample of 6 of the 25 pages
%  (seed 20260922: pp. 88, 89, 91, 101, 104; top-up seed 20260922-top-up:
%  p. 87) was read line by line at the image by readers who had not
%  transcribed them.  0 discrepancies (95% upper bound ~0.5/page).  The same
%  review re-adjudicated the batch agents' doubtful calls and overturned one
%  (p. 102 `religions-/filosofiske', joined -> hyphen kept, by the rule for
%  line-end hyphens); four italic calls spot-checked, all confirmed.
% ============================================================

\documentclass[12pt,a4paper,openany]{book}
\usepackage[utf8]{inputenc}
\usepackage[T1]{fontenc}
\usepackage{libertinus}
\usepackage{libertinust1math}
\usepackage{graphicx}    % for \reflectbox (the old Danish ɔ: id-est mark)
\DeclareUnicodeCharacter{0254}{\reflectbox{c}}  % ɔ (U+0254) = reversed c
\usepackage[danish]{babel}
\usepackage[protrusion=true,expansion=true]{microtype}
\usepackage{setspace}
\usepackage[margin=1.2in]{geometry}
\usepackage{fancyhdr}
\setlength{\headheight}{28pt}

% Footnotes as the 1903 printing sets them: *), not 1, 2, 3.  Set the mark
% explicitly at each site with \fnmark{*)}{...} (or \fnmark{**)}{...} should a
% page carry a second note); the default is restored afterwards.
\renewcommand{\thefootnote}{*)}
\newcommand{\fnmark}[2]{\renewcommand{\thefootnote}{#1}\footnote{#2}%
  \renewcommand{\thefootnote}{*)}}

% A piece's head.  #1 = the title as the volume's INDHOLD prints it (used for
% the table of contents and the running head); #2 = the head block as printed
% over the piece, with its own \opage.  The two readings differ, and both are
% kept (TRANSCRIPTION-PLAYBOOK §6).  The short rule under every head is
% printed, and is set by the macro.
\newcommand{\piecehead}[2]{\clearpage\phantomsection
  \addcontentsline{toc}{section}{#1}\markboth{#1}{#1}%
  \begin{center}\large #2\\[6pt]\rule{2cm}{0.4pt}\end{center}\par\medskip}
% The short centred rule printed at the end of every piece.
\newcommand{\piecerule}{\par\begin{center}\rule{1.5cm}{0.4pt}\end{center}\par}

\usepackage{hyperref}
\hypersetup{pdftitle={Brandes, Samlede Skrifter XIII: Anmeldelser og Indlæg om Rasmus Nielsen (1867)}, pdfauthor={Georg Brandes}, hidelinks}

\pagestyle{fancy}
\fancyhf{}
\fancyhead[LE,RO]{\thepage}
\fancyhead[RE]{\textit{Samlede Skrifter XIII (1903)}}
\fancyhead[LO]{\textit{\leftmark}}
\setstretch{1.3}

\title{\textbf{Anmeldelser og Indlæg om Rasmus Nielsen (1867)}\\[8pt]\large Georg Brandes, \textit{Samlede Skrifter}, Trettende Bind (Supplementbind),\\ pp.~85--109}
\author{}
\date{Kjøbenhavn: Gyldendalske Boghandels Forlag (F.\ Hegel \& Søn), 1903}

\usepackage{marginnote}
\newcommand{\opage}[1]{\marginnote{\footnotesize\textit{[#1]}}}
\newcommand{\apage}[1]{\marginnote{\footnotesize\textit{[#1]}}}
\begin{document}
\maketitle
\thispagestyle{empty}
\tableofcontents

\mainmatter

% ============================================================
%  Om den gode Vilje som Magt i Videnskaben (printed pp. 85--92)
%  Indhold (PDF 9) reads: `Om den gode Vilje som Magt i Videnskaben 85'
%  Head as printed: `R. Nielsen: Om den gode Vilje som Magt i' /
%    `Videnskaben' / `(1867)'  -- two italic lines, confirmed at the image.
%  Running head (pp. 86--92) reads: `Om den gode Vilje som Magt i
%    Videnskaben' (= Indhold).
%  The Nielsen quotation on pp. 88 and 90 is set as smaller block matter
%  ({\small ...}), each opening with a paragraph indent.
%  Printer's defects logged at the site: p. 87 stray comma `sig,';
%  p. 90 comma for full stop before `Idet'; p. 91 unclosed quotation
%  (‹ count +1 for this piece); p. 92 wrong-fount `o' in `Tro'.
% ============================================================
% --- p. 85 ---
\piecehead{Om den gode Vilje som Magt i Videnskaben}{\opage{85}\emph{R. Nielsen: Om den gode Vilje som Magt i}\\
\emph{Videnskaben}\\[2pt](1867)}
Som Svar paa Biskop Martensens Skrift: \emph{Om Tro og Viden}
har Professor \emph{Nielsen} paa den \emph{Gyldendalske Boghandel}s Forlag
udgivet \emph{Om den gode Vilje som Magt i Videnskaben}, en bred
og drøj, en haanende og bidende lille Bog, skreven med over\-%
strømmende Frembringelseslyst og fuld af gode Indfald. Dens
Lystighed har dog paa sine Steder en mere hollandsk end dansk
Karakter.

At Titlen er ironisk, da Bogen vil vise, at i Videnskaben
betyder den gode Vilje Intet, er en Oplysning, som det er nød\-%
vendigt at fremsætte for dem, der endnu erindrer, at der var en
Tid, da Professor Nielsen selv gjorde den teologiserede Samvittig\-%
hed til Udgangspunkt for den sande Videnskab, og da det fra
Professorens egne Læber forkyndtes, at f. Eks. Fritænkerens
Bibelkritik, skønt \emph{upersonlig sand}, var uden Værd, fordi den
havde sin Grund i en \emph{personlig Usandhed}. Det er saa meget
mere nødvendigt at fremhæve dette, som fossile Levninger af
denne Anskuelse endnu den Dag idag forekommer i de Ydelser,
der frembringes af Professor Nielsens nuværende Tilhængere,
saa disse ser, at de maa skynde sig noget, hvis de vil blive ved
at følge med og ikke som deres Forgængere sejles agterud og
tabes paa Vejen.

Bogen kan ses fra en dobbelt Side: den vil nedbryde og
den vil opføre. Den vil \emph{nedbryde} Teologien. Hvad er Teolo%
% --- p. 86 ---
\opage{86}gien? Det er en Disciplin, som staar i samme Forhold til Filo\-%
sofien som Alkymien til den videnskabelige Kemi og Astrologien
til Astronomien. Den har overlevet sine Aandsbeslægtede, den
er en af de mange Efterladenskaber fra Middelalderen, der lang\-%
somt smuldrer hen i det moderne Samfund. En Kamp mod denne
Døende har tabt enhver Interesse, den er et for den Dannede,
den Arbejdende, den Stræbende uværdigt Skuespil, hvorfra han
allerhelst vender Øjnene bort. Hvor kan Professor Nielsen da
vedblive at betragte Teologien som egnet til et Angreb med
skarpe Vaaben? Hvem angriber en værgeløs gammel Kvinde,
der neppe kan holde sig oprejst paa sine Krykker? Hvem sprin\-%
ger med dragen Klinge ombord paa en synkefærdig Skude, der
langsomt gaar tilbunds med Mand og med Mus?

Jeg véd ikke andet Svar end: Den, man elsker, tugter man.
Professor Nielsen svarer Biskop Martensen, som han svarede
Pastor Zeuthen; han har altid havt en Smule --- idetmindste en
Smule Trykkersværte --- tilovers for de Ordinerede, han sige,
hvad han vil. Det Eneste, der kunde undre, er det, at Pro\-%
fessor Nielsen i sin Opdrager-Iver har glemt den gamle smukke
Lærer-Regel: Hvis jeg ikke var hidsig, skulde Du faa Hug; ti
Professor Nielsen er hidsig, og han hugger ligefuldt, og derfor
hugger han vildt og slaar engang imellem sig selv over Fingrene
af lutter Iver.

En Læser, der har modtaget et samlet Indtryk af den nye
Tids kritiske Forskning, kan ikke Andet end beklage, at Pro\-%
fessor Nielsen under Indflydelse af en Heftighed, der er yngre
end hans Aar, lader Geologien anvise Paradiset Plads i ‹Amme\-%
stuen› og Botanikeren sende dets Træer ‹ad Helvede til›; ti
slige Udtryk vil ubehageligt minde ham om den Tid, da Viden\-%
skaben var polemisk, fordi den ikke endnu formaaede at være
kritisk, om den Tid, da man spottede Overlevering og Under
istedetfor at forstaa og forklare den sjælelige Oprindelse og den
historiske Udvikling af Under og Overlevering. Denne Tid er
nemlig forbi, og det kan ikke være Professor Nielsens Hensigt
at kalde den tilbage. Dertil har han dels en altfor grundig
Dannelse, dels en altfor grundig Velvilje for selve denne Teologi,
hvorfra han er gaaet ud, og til hvilken han bestandig uimod\-%
staaeligt føler sig dragen. Dette er et Punkt, man altid
maa holde fast, ellers vil man hvert Øjeblik føle sig forvirret.
Der findes f. Eks. i Professorens Bog Steder, hvoraf man ellers
% --- p. 87 ---
\opage{87}vilde kunne lade sig forlede til den urigtige Antagelse, at Pro\-%
fessorens sidste Ord egenlig faldt sammen med Fritænkerens
første, saaledes f. Eks. det udmærket vittige Stykke, der ender
med de Ord: ‹De (Teologerne) erkendte, at de ikke fremdeles
kunde vedblive at tale højrøstet om Historiens overnaturlige
Kendsgerninger, efterdi de overnaturlige Kendsgerninger nok ikke
var historiske og de historiske ikke overnaturlige›; men Ingen
maa lade sig forvirre af Sligt. Naar Professor Nielsen her be\-%
raaber sig paa Erfaringen og paa den ‹sunde Menneskeforstand›,
der som Slange giver Teologien at æde af Kundskabens Træ,
da mener han det ikke alvorligt, hans eget Standpunkt er ikke
‹den sunde Menneskeforstands›, men et ganske andet, som nu
skal blive paavist, idet vi gaar over til, hvad Professor Nielsen
vil \emph{opføre}.

‹Den sunde Menneskeforstand› indser f. Eks., at der er
netop ligesaa megen og netop ligesaa liden Grund til at antage
Dæmoner som til at antage Centaurer. En videnskabelig Kritiker
forklarer, hvorledes hos Oldtidens Folk Troen paa Dæmoner er
opstaaet. Professor Nielsen endelig bærer sig, saaledes ad: Han
% `sig,': a stray comma (or comma-shaped blot) at the baseline after
% `sig' -- transcribed as printed; both witnesses read it too.
begynder med at overhøre den sunde Menneskeforstands Ind\-%
sigelse; dernæst giver han Kritikeren Ret ‹fra Videnskabens
Standpunkt›; men føjer saa selv til dennes Videnskritik en
‹Troskritik›, der lyder egne Love og Intet har med Videns\-%
kritiken at gøre, ikke krydses af den, ikke støder sammen med
den hverken krigersk eller fredeligt, fordi ethvert Sammenstød
er umuligt, da de to Kritiker lyder hver et absolut forskelligt
Princip. Videnskritiken har afskaffet Dæmoner, men samtidig
forklaret dem; Troskritiken beholder dem overalt, hvor deres
Eksistens ‹anskueliggør Troens Grundmirakel›; at det hele Spørgs\-%
maal om Dæmoner for sund Menneskeforstand er uden al Inter\-%
esse og forud afgjort, kommer ikke i Betragtning. Hvorledes
gaar det nu til, naar Troskritiken faar at vide, hvor overlegent
Videnskritiken, der har sine Papirer i en anden Hjerneskuffe,
behandler Troskritikens egne Dæmoner? Hemmeligheden er: den
faar det ikke at vide. Hvorledes det? ‹I det Øjeblik Tros\-%
kritiken faar det at vide, er den nemlig ikke mere Troskritik,
men Videnskritik.› Det er umuligt, at Floden Tajo kan løbe
ind i Portugal, ti i samme Øjeblik den løber ind i Portugal,
hedder den ikke mere Tajo, men Tejo.

Troskritiken og Videnskritiken kan da saa lidet hindre eller
% --- p. 88 ---
\opage{88}indskrænke hinanden som Tro og Viden selv, hvilke Professor
Nielsen med rolig Ubekymrethed om alle Indvendinger vedbliver
at betegne som ‹absolut uensartede Principer›. Den nærmere
Udvikling heraf er følgende, højst besynderlige:

% block quotation set in smaller type, opening with an indent
{\small
Bevidstheden halveres ved relativt ensartede, men ikke ved absolut
uensartede Faktorer. Naar en Bevidsthed halveres i Tanken, faar den en
halv Tanke istedetfor en hel: det, der mangler, er Tanke; halveres den i
Vilje, faar den en halv Vilje istedetfor en hel: det, der mangler, er Vilje.
Derimod kan man ikke i strengere Forstand sige, at Bevidstheden halveres
af Tanken og Viljen, saaledes at Tankens Halvdel bliver Vilje, Viljens
Halvdel Tanke [kan man da sige dette Nonsens i nogensomhelst For\-%
stand?], ti Tanke kan kun i højst uegenlig Mening forvandles til Vilje og
Vilje til Tanke. [Det er naturligt nok, at det uegenlige Udtryk --- For\-%
vandling --- maa tages i uegenlig og ikke i egenlig Forstand]. Staar det
da fast, at Troen paa ingen Maade [nylig var det kun: i egenlig Mening
ikke] lader sig forvandle til Viden eller Viden til Tro, saa staar det ogsaa
fast, at Troen ligesaa umulig kan halveres ved Viden som Viden ved Tro.
Kan en Bevidsthed ikke halveres af Venskab og Matematik, af Kærlighed
og Kemi, saa kan den endnu mindre halveres af Etik og Astronomi, af
Religion og Geologi.\par}

Det er en Ulempe, at der her netop er forudsat det, som
skulde bevises. At en sand Moral og en sand Religiøsitet ikke
kan støde an imod en sand Videnskabelighed, forstaar sig af
sig selv. Men Spørgsmaalet er netop, hvilken Moral og hvilken
Religiøsitet der er den sande, en positiv eller en fri? Spørgs\-%
maalet er f. Eks., om en lille, med sig selv tidt nok uenig,
semitisk Stammes religiøse Oldtidsforestillinger fra de forskel\-%
ligste Tider, Forestillinger, der paa alle Punkter er gennem\-%
vævede med og baarne af dens hele Verdensanskuelse, dens
Tanker om Jordens Tilblivelse, dens Overtro, dens utallige For\-%
domme, om disse modsiger eller ikke modsiger vor Geologi og
vor Astronomi. De modsiger dem, paa samme Maade som de
i utallige Tilfælde modsiger vor Moral og vor religiøse Følelse.
Men lykkes det nu paa ingen Maade Professor Nielsen at holde
de to Halvkugler ude fra hinanden, saa lykkes det ham dog
endnu langt mindre til Slutning at forene dem. Ti saa skarp
end Adskillelsen paastaas at være, saa maa jo dog en Forening
kunne tænkes: ‹De to absolut uensartede Principer mødes,
skilles og træder \emph{paa en Maade} i Underhandling med hinanden i
\emph{Meningens} Medium›. Meningen --- det er den frelsende Kategori;
paa Meningens neutrale Grund underhandler de to absolut Uens%
% --- p. 89 ---
\opage{89}artede med hinanden. Den Eneste, der kunde have Noget mod
en saadan Underhandling, det skulde da være det danske Sprog;
ti vel er det Filosofernes Forret at lære Sproget Ræson;
men de to Uensartede forholder sig jo som Kærlighed og Mate\-%
matik; hvad hedder det Medium, hvor Kærlighed og Matematik
forhandler sammen?

Professor Nielsen oplyser ved et Eksempel: ‹Hvad der i første
Mosebog fortælles om Paradiset, om Livets Træ, Kundskabens
Træ osv., er ganske vist fortalt i den Mening, at det Altsammen
hører under det udvortes Virkelige, at Livets Træ og Kund\-%
skabens Træ ligesaa vist har deres Rødder i Havens sanselige
Jordbund som Vintræet eller Figentræet. Men i denne Mening
kan vi \emph{paa Kulturbevidsthedens nuværende Trin} umulig opfatte
det?›

Læste vi ret? Kulturbevidstheden? Glemmer Professor Nielsen
sig selv og sine Paastande? Hvad har Kulturbevidstheden vel
her at gøre? Skal Troskritiken grundes paa Kulturbevidstheden,
og Troen endda være absolut uensartet med denne, ikke til at
bestemme og ikke til at indskrænke ud fra den? Hvilket Hjerne\-%
spind, hvilket Virvar, ud af hvilket der ingen Udvej gives!
Moral: Ræk ikke Kulturbevidstheden den lille Finger, ellers
tager den hurtigt den hele Haand.

‹Hine Billeder,› siger Professor Nielsen, ‹har ikke ydre men
indre Virkelighed.› At dette ikke er nyt, det véd Enhver, der
kender noget til Teologiens Historie; at det er uforstaaeligt, det
erfarer Enhver, der forsøger paa at tænke det igennem. Hvad
er indre Virkelighed, naar den skal være forskellig fra poetisk?
--- ‹Sagnet om Paradis,› siger Professor Nielsen, ‹kan, saa sandt
det i rene, klare Træk for Troens Blik afslører et af de dybeste
Viljesmysterier, umuligt være Myte.› Gør Myten om Prome\-%
theus maaske ikke det? Og er dette et Skelnemærke? Naar her
fraskrives Sagnet udvortes Virkelighed, er det ligegyldigt, om
det kaldes Myte eller ej. At Paradiset har en poetisk Virkelig\-%
hed, som ‹Jura› og ‹Graavakke› ikke har, forstaar sig af sig
selv og negtes af Ingen.

Til den sjælelige Forening i \emph{Meningens} Himmelegn svarer
nu den væsenlige (ontologiske) i \emph{Mysteriets}. Udviklingen af dette
Punkt er given i et Sprog, der beklageligt minder om den Hegel\-%
ske Tids uforstaaeligste Mundart.

% block quotation set in smaller type, opening with an indent
{\small
% --- p. 90 ---
\opage{90}Tro og Viden lader sig kun forene i et absolut Mysterium. Det
absolute Mysterium er paa éngang eksklusiv Grænse og negativ Midte.
Som eksklusiv Grænse sætter Mysteriet Skilsmisse mellem Tro og Viden
som imellem to Regioner. Naar den Troende ser mod Troens, den Vi\-%
dende mod Videns Midtpunkt, bliver det, der er sandt i Troen, usandt i
Viden og omvendt. Som negativ Midte forener Mysteriet derimod de
adskilte Ekstremer og viser sig som Regionernes Gennemgangsmedium,
% `Gennemgangsmedium,' before capital `Idet': a comma where the sense
% wants a full stop -- confirmed a comma at 600 dpi, transcribed as printed.
Idet den Troende og den Vidende begge paa én Gang retter deres Blik
mod Mysteriet, udlignes Striden, og Ekstremerne enes. Set fra Videns
Standpunkt er nu Miraklet muligt [ja, hvad gør ikke en punktlig Ekser\-%
% `gør': the stroke of the ø is not visible in this copy (it reads `gor'
% at 600 dpi); both witnesses read `gør', and the stroke of ø is faint
% throughout this small type (cf. `tilhøjre' below). Taken as ø.
cits, hvor der paa én Gang ses tilhøjre og tilvenstre!]. Ti vel er det
umuligt, at der nogensinde skulde kunne ske et Brud paa de evige, med
Fornuftnødvendigheden identiske Naturlove; men med Lovenes Uforander\-%
lighed kan jo, naar de i Mysteriet skjulte Mellembestemmelser under\-%
forstaas, den Aabenbarelse af Viljens, Visdommens og Kærlighedens Magt,
hvori den Troende ser (!) Miraklet, meget vel tænkes at bestaa. Set fra
Troens Standpunkt kan Fornuftnødvendigheden med de uforanderlige
Naturlove forenes med Forudsætningen af Miraklet. Ti vel er det umu\-%
ligt, at Nødvendigheden skulde kunne hæmme Friheden, og Naturens
Love indskrænke Guds Vilje; men naar de i Mysteriet skjulte Mellem\-%
bestemmelser underforstaas, lader det sig jo tænke, at Naturlovene, langt
fra at brydes, tvertimod i allerhøjeste Forstand opfyldes, hver Gang en
guddommelig Viljesakt fuldbyrdes.\par}

Dette var i gamle Dage god Teologi. Filosofiske Antyd\-%
ninger af noget Saadant finder man f. Eks. hos den tyske Præst
Bockshammer, der anføres af J. L. Heiberg i \emph{Prosaiske Skrifter} I, 110.
Denne slutter med en flydende Ordstrøm, der indeholder selv\-%
samme Paastand. Ti mere end en ubevist og ubevislig Paa\-%
stand er dette jo ikke. Lad os vælge et Under af det nye
Testamente, f. Eks. Figentræets Forkullen ved det guddomme\-%
lige Magtord, der straffer det for ikke at bære Frugter. Man
behøver dog ikke at være Tænker for at indse, at et Mirakel
som dette, hvis sjælelige Oprindelse iøvrigt synes at være Mis\-%
forstaaelse af en Parabel, som om denne var historisk Virke\-%
lighed, at et saadant Mirakel ikke lader sig bringe i Overens\-%
stemmelse med Naturens Love hverken i jævn og simpel For\-%
stand eller i nogen højere eller ‹allerhøjeste› Forstand. Aller\-%
højstsamme Forstand vil her kun spilde sine Kræfter.

Tvekampen mellem Biskop Martensen og Professor Nielsen,
disse to saa øvede og udviklede Kæmpere, gør --- det være
sagt med oprigtig Beklagelse --- samme Indtryk som en Øvelses\-%
strid imellem to lige slette Fægtere: Ingen af dem kan bruge sin
Huggert; Ingen af dem afbøder Modstanderens Hug; men det
% --- p. 91 ---
\opage{91}gælder for dem begge kun om, hvem der kan bibringe den
anden de hurtigste og drøjeste Slag. De slaar som de gamle
nordiske Bersærker med Skjoldene paa Ryggen, og man faar ikke
Lejlighed til at glæde sig over stor Kunst eller Behændighed
paa nogen af Siderne. Men er her ikke megen Kunst, saa er
her desto mere Politik.

Det var god Politik af Hans Højærværdighed, Biskop Mar\-%
tensen at lade Professor Nielsens Angreb paa hans Dogmatik
uden Svar; ti det gaar med Dogmatiker nutildags som med
andre Fæstninger, der for hundrede Aar siden syntes omtrent
uindtagelige: de lader sig ikke mere forsvare; det gaar ned ad
Bakke med deres Volde, deres Grave er Rendestene, og vil
Fæstningen ikke udlevere sine Nøgler, gør det ikke Noget, efter\-%
som der ingen Porte er.

Men ligesaa sikkert var det god Politik af den filosofiske
Kæmper at gaa angrebsvis til Værks og besvare Biskop Marten\-%
sens høflige ‹Goddag› med et frygteligt ‹Økseskaft›. Professor
Nielsens Hovedvagt duer ikke til Forsvar. Professor Nielsen
har vel én Mand i Vagten, men han græder --- med det ene
Øje og ler med det andet; hans højre Haand véd ikke, hvad
hans venstre gør. Han bevæger sig frem paa to uensartede
Ben, et Kødben og et Træben, paa to absolut spaltede Ben,
flækket som den ene af Uffes Saksere. Hans Hoved er et Janus\-%
hoved: med den ene Mund forkynder han Naturlovens højere
Opfyldelse igennem Miraklet; med sin anden Strube giver han
Kulturbevidstheden Æren og synger en velment Lovsang for
David Strauss. Dette sidste er maaske det særeste blandt mange
sære Træk i denne Fejde. Menigheden paamindes af Professor
Nielsen om langt fra at tage Forargelse af Strauss's Forvovenhed,
at paaskønne den Hjælp, som denne Mand ved sin ‹konsekvente
og resolute› Fremgangsmaade umiddelbart har ydet Troen, idet
han har grundlagt den Videnskritik, hvilken Professor Nielsen
ved sin planlagte Troskritik vil give det nødvendige Sidestykke.
Biskop Martensen maa høre ilde for sin rettroende Fordøm\-%
melse af Strauss. ‹Skal Strauss ekskommuniceres, da er der
Mange, som maa ekskommuniceres [ja, hvorfor ikke, det er jo
% the quotation opened at `‹Skal' is never closed: no › before `[ja'
% (checked at 600 dpi) -- printer's omission, transcribed as printed.
en god kristelig Lære!]; ‹er Strauss en Antikrist, da er der i
Lejren mange Antikrister› [ja, hvilken Ortodoks tvivler vel paa
det!] Janushovedet henvender sig til begge Sider, det staar
som Don Juan imellem Bønderpigerne og har fagre Ord for
% --- p. 92 ---
\opage{92}begge Parter. Grundtvigianerne faar deres Del af disse. \emph{Dansk
Kirketidende} modtager offenlig Ros, Hr. N. Lindberg prises ---
hvorfor tror man? --- for \emph{logisk-kategorisk Sikkerhed}.

Narren siger i \emph{Aladdin}: ‹Jeg gantes ej, det staar min Tro
% `Tro': the o is a wrong-fount (heavier, rounder) letter.
i Brevet.›

Hvis nogen vil gøre sig den Ulejlighed at opsøge den
Artikel, hvorom Talen er, da vil han, efter Sigende, i dens
Noter kunne finde interessante Prøver paa den Sikkerhed, som
Professor Nielsen fremhæver, og saaledes med egne Øjne kunne
overbevise sig ikke om den Maade, hvorpaa en Kirketidende
deltager i en videnskabelig Polemik, ti hvem søger Videnskab i
Kirketidender, men om hvad i Danmark en Filosof kan rose.
Professor Nielsen glemmer ikke sine Pappenheimere; hans Pap\-%
penheimere vil heller ikke glemme ham.
\piecerule

% ============================================================
%  Nielsen og Grundtvig (printed pp. 93--97)
%  Indhold (PDF 9) reads: `Nielsen og Grundtvig 93'
%  Head as printed: `Gb. (Rudolf Schmidt):' / `R. Nielsens Filosofi og den
%    Grundtvigske Anskuelse.' / `(1867).'  -- two italic lines (the
%    parentheses of line 1 italic with it), then the year in roman WITH a
%    full stop after the parenthesis, unlike the other heads.  Short rule.
%  Running head (inner pages) reads: `Nielsen og Grundtvig' (= Indhold).
%  Unsigned, undated; ends with a short rule.  p. 97 foot carries the
%    signature line `G. Brandes: Samlede Skrifter. XIII.' and `7' -- not
%    transcribed.
%  `Gb.' is italic where it names the pseudonym (pp. 93, 97).
% ============================================================
% --- p. 93 ---
\piecehead{Nielsen og Grundtvig}{\opage{93}\emph{Gb. (Rudolf Schmidt):}\\
\emph{R. Nielsens Filosofi og den Grundtvigske Anskuelse.}\\[2pt](1867).}
Under Mærket \emph{Gb.} har en Forfatter, der staar Grundtvig og
Nielsen lige nær, udgivet en Bog paa 124 Sider, betitlet \emph{R. Niel\-%
sens Filosofi og den Grundtvigske Anskuelse}, en Bog, hvis Form
vidner om megen stilistisk Modenhed og Dannelse. Dens Stand\-%
punkt er det, at give Nielsen ubetinget Ret, baade hvor han an\-%
griber, og hvor han er Genstand for Angreb, at hævde Ægthe\-%
den af Grundtvigs Genialitet og Betydningen af hans Person og
Lære, men samtidig kritisk at udskille en Del af det meget Løse
og Frastødende, der hænger sammen med Grundtvigs hele Virk\-%
somhed, endelig til Slutning, saa godt det lader sig gøre, at
rime Nielsen og Grundtvig sammen. Bogens Indholdsværd er
et meget forskelligt i de forskellige Partier. Kortest og fyndigst
lader det sig bestemme saaledes, at det større Parti, Afsnittet
om Grundtvig, der er forfattet med en agtværdig Stræben efter
Upartiskhed, og som hviler paa Grundlag af en ualmindelig Be\-%
læsthed i dansk Literatur, indeholder meget Sandt og Træffende,
men at derimod de to korte Afsnit om Nielsen, der begynder
og slutter Bogen, er ligefrem usle. De er ikke blot tomme for
Tanker, men skrevne paa en næsvis og uanstændig Maade. Ti
det er Næsvished at indblande sig i en alvorlig Strid for med
overlegen Fornemhed at afgive sin uforgribelige Mening som en
Højesteretsdom, naar man ikke formaar at levere et eneste Bi\-%
drag til Striden, eller at gendrive en eneste af de Bevisgrunde,
% --- p. 94 ---
\opage{94}man slaar paa Nakken af, men nøjes med den Floskel, at An\-%
griberne har vist sig at være ‹uden Tilegnelsens Kraft overfor
det Nye›.

Alt, hvad der indtil Dato er forebragt imod den nye Lære,
affærdiges med den Snak, at man dog umuligt, naar man aner\-%
kender Nielsens ualmindelige Begavelse, kan tiltro ham saadanne
Fejl og Modsigelser som dem, der tillægges ham. En større
Bagvendthed end den, der ligger til Grund for dette ofte nok
gentagne Svar, lader sig ikke tænke. Sæt, en Mand skal gen\-%
nemgaa et vanskeligt Regnskab. Dette Gennemsyn anses for en
Formsag! ti Omgivelserne betragter det som en Trosartikel, at
dette Regnskab skal være rigtigt; saaledes befandtes det altid af
de tidligere Revisorer, og gjorde disse Vanskelighed, fik man
Sagen dysset ned. Manden vilde komme i Strid med sin egen
Tillid til Autoriteterne og lægge sig ud med alle pæne Folk, i
Fald han ikke fik det givne Facit ud; han faar det da ogsaa
lykkeligt til at stemme. Uheldigvis opdager nu En, at dette hel\-%
dige Udfald skyldes den Omstændighed, at et Steds i Regnska\-%
bet ved en Fejltagelse 2 og 2 er lagte sammen til 5. Han gør
opmærksom derpaa og faar et efter det andet fire Svar: \emph{a}) Denne
Indvending beror paa en Misforstaaelse. \emph{b}) Du maa da kunne
sige Dig selv, at en Regnemester som denne umuligt \emph{kan} be\-%
gaa en saadan Fejl. \emph{c}) Hvad Du siger er afskyeligt og oprø\-%
rende, undergraver al Fortrøstning til det Bestaaende og hele
den barnlige Tillid til Autoriteterne, uden hvilken Alt gaar fra
hinanden. \emph{d}) Hvilken forslidt og fortærsket Indvending! I
Sandhed en højst ny og mærkelig Opdagelse den, at 2 og 2
er 4!
% the list letters a)--d) are italic, their closing parentheses roman
%   (upright), here and below.

\emph{a}) Indvendingerne beror paa Misforstaaelse. Svaret er sim\-%
pelt. Hvor længe tror I, at man vil taale denne Udflugt? Mis\-%
forstaar vi, saa luk Munden op og tal tydeligt, saa man kan for\-%
staa Jer, i Stedet for at tale, som I havde Grød i Munden!
Det er kommet saa vidt, at det nu er blevet Mode i Filosofien,
som det længe har været i Politiken, at enhver Udtalelses høje\-%
ste Kunst er den: at skjule, hvor man egenlig vil hen, at ud\-%
trykke sig i Gaader, at sige Ja og Nej som en diplomatisk De\-%
peche eller en Tale af Napoleon III.

\emph{b}) Skulde Professor Nielsen ikke vide, at 2 og 2 er 4?

Dette angaar ikke os og er os ganske ligegyldigt; men vi
paaviser den Kendsgerning, at han her paa dette bestemte Punkt
% --- p. 95 ---
\opage{95}har ladet 5 være lige. Kan Noget tænkes latterligere, end at
man lige overfor den bestemte Paavisning vedbliver at beraabe
sig paa Mandens Evner. Der er paavist Modsigelser Punkt for
Punkt i Professor Nielsens Lære. Der gives intet andet Vaaben
imod Logik end Logik. Men vi spørger logisk, I svarer os ly\-%
risk. Det er, som om den berømte Mosekone var bleven den
hele Lejrs Valkyrje. Man kommer med Fremstillinger, der gaar
i Spring, med fantasirige Udmalinger, Lignelser, Skyggebilleder,
Alt undtagen Bevisførelse, Gendrivelse, Logik; ti af dette findes
her end ikke Skygge.

\emph{c}) Indvendingerne vil Idealerne til Livs (se Anmeldelsen af
denne Bog i \emph{Fædrelandet}), vil med tør Forstandighed røve
% the `)' after Fædrelandet looks slanted; not clearly italic, kept outside.
Troen dens Grund og Begejstringen dens Kilde. Vi svarer: Blæse
med slige Idealer, til Hekkenfeldt med en saadan Begejstring,
der kryber i Skjul i \emph{Eksistensens} Smuthul, naar Videnskaben ud\-%
æsker i Aandens Sollys. Hold Jer til Fakta og gendriv vore
Bevisgrunde! Lad hine store Ord: Idealerne, Begejstringen, der
opslaas som Plakater og imponerer Hoben, fare og stiller Bevis
imod Bevis. Ros Jer dog ikke af Jer nøjagtige Metode, der
skal gøre Eders Slutningsrække saa paalidelig som den matema\-%
tiske, naar Alt, hvad I siger, har en blot skønliterær og ingensom\-%
helst videnskabelig Karakter. De skønliterære Sandheders Tid er
forbi. Ve den, der ikke kan taale den nøgne Sandhed, ve den,
der sætter Filosofiens Opgave i at klistre Figenbladet paa den!

\emph{d}) Hvad der forebringes imod Professor Nielsen, er gam\-%
melt. Ja, det er gammelt, at en Tænker ikke maa modsige sig
selv; ja, Sandheden er gammel, gammel som alt Evigt, men
hvilken Anke at komme med den! Er der Noget, som er usselt
i den danske Ungdom, saa er det dens Tilbøjelighed til Slutnin\-%
ger som disse: ‹Dette er nyt, altsaa sandt; dette er nyt, altsaa
vort; dette er nyt, altsaa det, hvorom vi skal stimle‹. Nyheds\-%
% PRINTER'S ERROR: the quotation closes with a TURNED mark (opening
%   guillemet `‹' for `›') after `stimle'; transcribed as printed.
%   Guillemets on this page therefore run 5 opening : 3 closing.
kræmmerne hører ikke Videnskabens Tilraab: ‹I Daarer! prøv
dog først dette Nye!› ‹Nej,› hedder det, ‹lad os ile, lad os haste
for at komme dybt ind i Trængslen og presse den bort, der
gaar imod Strømmen!› Og tror man da, at Nogen lader sig af\-%
skrække deraf? Tror man da, at Nogen, der i slige Spørgsmaal,
saa vigtige, saa gennemgribende, saa kvalfulde, ene med sig selv,
men Ansigt til Ansigt med Videnskabens Aand har fravristet
denne en Overbevisning, der er levende og kraftig som Aanden
selv, bekymrer sig stort om disse Tidens Strømninger paa Over%
% --- p. 96 ---
\opage{96}fladen eller om dem, der kun føler sig trygge og stærke, naar
de midt i Trængslen, hvor Ingen kan røre Arm eller Ben, bæ\-%
res viljeløst frem i en eller anden Beundrets Kølvand? Vi har
ingen Klake i vore Teatre, skal vi nu faa en Klake i vor Li\-%
teratur?

Nej, det er umuligt at opgive Haabet om, at Ungdommen
i Danmark trods al dens Magelighed vil forsøge paa at naa til
selvstændigt Skøn. Lad den prøve sig frem, lad den fatte den
Beslutning engang at ville se med egne Øjne! Lykkeligvis er vi
jo ikke aandeligt stavnsbundne; der er Udveje for den, der føler,
% `føler': the ø is damaged -- its bar barely shows and the letter reads
%   almost as `o' (witness 1 `foler'); read as ø, the sort being defective.
at vi i en Gren af vor Literatur staar tilbage for det øvrige
Europa. Lad Ungdommen studere fremmede Tænkere, og det
skulde i Sandhed gaa underligt til, om den ikke opdagede, at
disse Deklamationer ikke er Gendrivelser; lad den studere den
nyere Teisme f. Eks., og den vil se Sammenhængen imellem
Professor Nielsens Bestræbelser og den europæiske Reaktion
mod formentligt ensidig Panteisme; lad den blot faa Horisont,
vid Horisont, og den vil opdage, hvor lidet enestaaende Pro\-%
fessor Nielsens ‹udødelige› Behandling af Subjektivitetsproblemet
er, denne Behandling, der i Skriftet ses i Sammenhæng med den
oldnordiske Mytologi.

Hvad der her er virksomt til at vildlede Mange, det er føl\-%
gende Ræsonnement: Naar det maa indrømmes af Enhver, at
de religiøse Problemer (først og sidst af Kierkegaard) i vor egen
Literatur er behandlede i et Omfang og en Dybde, som maaske
neppe i nogen europæisk Bogverden, hvor kan man da tro, at
vi i nogen anden skulde kunne finde større Klarhed over disse
Spørgsmaal end hos os! Dette Ræsonnement er meget bestik\-%
kende, men skønt tilsyneladende kraftigt, er det ganske hult. Ti
det er ingenlunde altid Tilfældet, at den stærkeste Krise hidfører
det bedste og sandeste Resultat. Ingen kan benegte, at imod
den franske Revolution 1789 og de følgende franske Revolutio\-%
ner, disse Kæmpe-Anstrengelser for at tilkæmpe sig Frihed, er
vor Revolution 1848 --- det lange Tog af pæne Mænd, der med
hinanden under Armen spaserede til Christiansborg --- skønt den
i sig selv ikke mangler Værdighed, ynkelig, spag og for Intet
at regne. Og dog --- hvor er der vel nu mest Frihed, i Frank\-%
rig eller her? Men omvendt forholder det sig paa det religions-%
% `religions-/filosofiske' at a line end: hyphen KEPT as a compound hyphen,
%   the volume prints both forms mid-line (hyphenated pp. 97, 98, 100; `religionsfilosofisk' solid p. 108).
filosofiske Omraade. Hvor er Frankrigs religiøse Kriser i dette
Aarhundrede ringe imod vore! Og dog --- hvor er der vel mest
% --- p. 97 ---
\opage{97}Sandhed, i den franske religions-filosofiske Kritik eller i vore
Frembringelser i lignende Retning? Den der har studeret den
franske Literatur blot i noget Omfang, vil neppe tvivle paa,
hvem han paa dette Felt maa række Palmen. Ti vel er Kierke\-%
gaard --- og om Ingen uden ham kan her være Tale --- en
uhyre Aand, en af dem, som neppe fødes én Gang hvert Aar\-%
hundrede, men Kierkegaard er vor Filosofis Tycho Brahe. Han
er stor som Tycho Brahe, men ligesom Astronomen vover han
i sin blinde Ærefrygt for Autoriteten ikke at sætte vort Systems
Centrum i dets Sol. I Astronomien er denne Sol et Legeme; i
Filosofien kaldes den Fornuften. Det synes som en tragisk
Skæbne, knyttet til Danmark, at dets største Opdagere skal tage
fejl i det Centrale. Men saa lidet som Eftertiden ved Tycho
Brahes Fejl forhindredes fra at tage hans Værk til Indtægt, saa
lidet falder Kierkegaard med det Positive i Moral og Religion.
Dette lader sig udskille af hans Fremstilling, og endda vil uende\-%
lige Skatte blive tilbage. De, for hvem dette som for Hr. Pro\-%
fessor R. Nielsen eller Hr. \emph{Gb.} er en lukket Bog, de kæmper
kun imod Fremskridtet, og stræber, om end forgæves, at hindre
det sande Nye i at bryde sig sin Vej. Men dette Nye er uimod\-%
staaeligt; ti det fører Fornuft og Frihed i sit Skjold.
\piecerule

% ============================================================
%  Et Svar (printed pp. 98--101)
%  Indhold (PDF 9) reads: `Et Svar 98'
%  Head as printed: `Et Svar' / `(1867)'  -- one italic line; the year in
%    roman, no full stop.  Short rule.
%  Running head (inner pages) reads: `Et Svar'.
%  Unsigned, undated; ends with a short broken rule (p. 101).  p. 99 foot
%    carries the signature `7*' -- not transcribed.
%  In this piece `Gb' / `Gb.' is set in ROMAN (p. 98), unlike pp. 93, 97.
% ============================================================
% --- p. 98 ---
\piecehead{Et Svar}{\opage{98}\emph{Et Svar}\\[2pt](1867)}
Forfatteren af Skriftet \emph{R. Nielsens Filosofi og den Grundtvigske
Anskuelse} har i \emph{Berlingske Tidende} for igaar indrykket et Gen\-%
svar til mit Angreb paa hans Bog, hvilket uden at forebringe
Nogetsomhelst mod de fire Punkter, jeg har fremdraget, ender
med de Ord, at ‹de, som virkelig tror, at her er en Sag, ikke
vil tage fejl af, at han (ɔ: jeg), stukken af Selvvigtighedens Braad,
kun danser Sanct-Veitsdans uden om de store Opgaver›. Slige
Ting er det tilstrækkeligt at anføre, ufornødent at besvare.

Med Hensyn til de Beskyldninger og Spørgsmaal, som min
ærede Modstander forøvrigt opkaster, synes det mig imidlertid,
at jeg skylder ikke ham, men \emph{Dagbladets} Læsere et Svar. Han
bebrejder mig at have henvist til den franske religions-filosofiske
Kritik, ‹som Avislæsere ikke kender.› Jeg indser ikke, hvad
Ondt der var heri, selv om det forholdt sig saaledes; men det
gør det naturligvis ikke. I vore Dage udgør Avislæsere som
bekendt ikke en Kaste for sig; saagodtsom alle Dannede læser
Aviser. Min ærede Modstander er selv et Eksempel. Vistnok
er Hr. Gb en særdeles dannet Mand (se f. Eks. Prøven paa hans
% `Gb en': no full stop after `Gb' (checked at 600 dpi); as printed.
Stil ovenfor), men hvorfor skulde \emph{Dagbladet} ikke være saa lykke\-%
ligt iblandt sine Læsere at tælle flere ligesaa dannede Mænd,
Mænd, der blandt andre Forfattere kender Renan, som Hr. Gb.
jo kender saa godt?

Min Modstander beskylder mig dernæst for at have henvist
religiøse Forestillinger, ‹der har baaret Menneskeheden oppe
gennem atten Aarhundreder›, til Hekkenfeldt. Det vilde unegtelig
være et hæsligt Træk. Men er det nu ogsaa ganske vist, at jeg
% --- p. 99 ---
\opage{99}har handlet saaledes? Er det vist, at de Forestillinger, jeg vil
have forvist, i atten Aarhundreder har \emph{baaret Menneskene oppe?}
(og dette er jo dog langt mindre væsenligt, end om de er egnede
til at bære Menneskeheden oppe den Dag idag). Føljetonen, der
forleden ved et Tilfælde ledsagede min Artikel, Brudstykket om
Middelalderens Hekseprocesser, der netop fremavledes af de
overnaturlige Forestillinger, som jeg bestrider, afgiver et godt
Bidrag til dette Spørgsmaals Besvarelse. I hin Tid, da det
Overnaturlige virkelig var levende i Menneskenes Bevidsthed,
straffede man med Kætterbaalet de Anskuelser, som man nutil\-%
% `med': the e is raised/damaged in the scan; the word is `med'.
dags nøjes med at lægge for Had. --- Et Par Ord, jeg kunde
sige til Forsvar for Ligheden imellem de religiøse og de politi\-%
ske Kriser, vil jeg spare, da min Modstander har pustet sig saa\-%
ledes op, at han \emph{forbyder} mig \emph{at tale} derom.

Der staar da blot ét Punkt tilbage, det atter og atter gen\-%
tagne Spørgsmaal, om hvis knusende Kraft man er saa over\-%
bevist: Hvorledes jeg kan anerkende Kierkegaards Storhed og
dog være overbevist om, at han har taget fejl i det Centrale?
Jeg kunde besvare dette Spørgsmaal, som det før er blevet be\-%
svaret af mig, med de Ord, at kun meget uselvstændige Mennesker
ved beundrende Anerkendelse ikke forstaar Andet end simpel Ved\-%
hængen gennem Tykt og Tyndt; men siden dette Spørgsmaal
paany drages frem, fremsætter jeg endnu følgende korte Med\-%
delelse om, hvorledes jeg er kommen til dette Resultat: Jeg
opdagede temmelig tidligt, at Kierkegaard ikke var ufejlbar. Jeg
opdagede f. Eks., og jeg antager, at det er gaaet Adskillige som
mig, at Kierkegaard var Tilhænger af den politiske Enevælde i
Danmark, at han paa det politiske Omraade med en mærkværdig
Indskrænkethed i Anerkendelse af det Nedarvede blot som ned\-%
arvet misforstod sin Tid og kæmpede imod Udvikling og Frem\-%
skridt. Jeg saa, at hans politiske Spaadomme, hvoraf han selv
er saa stolt, aldeles ikke slog ind. Jeg lededes saaledes til den
Iagttagelse, at han var Absolutist ogsaa paa det religiøse Omraade,
at det at bryde med Overleveringen, at svigte Forældres, at for\-%
lade en Faders Synsmaade var ham ikke blot en Rædsel, som
den vel føles, om end i forskellig Grad, af Enhver, der tvinges
fil at se Pieteten under Øjne, altsaa ikke blot en Angst, som
% PRINTER'S ERROR: `fil' for `til' (both witnesses and the image agree);
%   as printed.
det er Menneskets Opgave at overvinde, men en \emph{aandelig} Rædsel
at samme Art, som den, Naturen indplanter for at forhindre
% PRINTER'S ERROR: `at samme Art' for `af samme Art'; as printed.
Mennesket i at begaa en Brøde imod Naturen. Jeg antog og
% --- p. 100 ---
\opage{100}antager, at han har Uret heri; jeg holder Intet for helligt blot
som overleveret; ti at det er overleveret, gør det hverken for\-%
nuftigt eller sandt. Jeg betragter det som utvivlsomt, at man
angaaende dette Punkt kan lære mange Gange Mere af Tysk\-%
lands og Frankrigs religions-filosofiske Skribenter end af ham.

Det er bekendt, at Kierkegaards Værker udmærker sig ved
en sjælden Følgerigtighed og Enhed. Men det er, om maaske
end ikke almindelig anerkendt, saa dog neppe derfor mindre
afgjort, at det Forsøg, han paa et fremrykket Stadium af sin
Forfatterbane gjorde paa at opvise den inderlige Sammenhæng
imellem alle sine Skrifter dels er meget tvungent, dels ikke kan
udelukke, at disse Værker beholder deres Værd i og for sig uden
Hensyn til deres Stilling i Skrifternes System. Er Kierkegaard
vel paa dette Punkt Andet end det, han allernødigst vilde være,
og som han hadede allermest --- en Systematiker, der bringer
den frie Livsfylde i en forudfattet eller rettere tildels bagefter
fattet Grundtankes Snørebaand? Selv er han smittet af det syste\-%
matiske Raseri, han bekæmpede hos sin Tid. Men om det
endog var anderledes, om endog Alt, hvad Kierkegaard har
skrevet, virkelig fra først til sidst gik ud paa Fremstillingen
af det Særlig-Kristelige i den Form, hvori han har opfattet
det, vilde det derfor være nogen Anderledessindet forment at
glæde sig ved og tilegne sig de talrige Frembringelser af ham,
der Intet har med den paradokse Religiøsitet at gøre? Bliver,
set fra et fremmed Synspunkt, Afhandlingen om Don Juan min\-%
dre beaandet, \emph{In vino veritas} mindre aandfuld, Afhandlingerne
om Ægteskabet mindre fængslende eller Fremstillingen af den
ikke-kristelige Religiøsitet (i \emph{Afsluttende Efterskrift}) mindre betyd\-%
ningsfuld og mindre lærerig? Hvilken Urimelighed!

Men man behøver slet ikke at blive staaende herved. I
mangfoldige Tilfælde, hvor Kierkegaard efter sin egen Mening
fremstiller det Paradokse (som i \emph{Kærlighedens Gerninger}) er dette
Paradokse slet ikke fornuftstridigt og lader sig oversætte paa
det Humanes Sprog; i andre Tilfælde, f. Eks. i \emph{Sygdommen
til Døden}, kommer det Positive henimod Slutningen uafledet ind.
Det er umuligt ud fra Fortvivlelsens Begreb, som Kierkegaard
i Begyndelsen rent menneskeligt fastsætter det, at faa Fortviv\-%
lelsen og dermed Synden til at naa Højdepunktet i Rationalis\-%
mens Betragtning af det Kristelige. Heri er da Meget, som en
% --- p. 101 ---
\opage{101}Fornuftdyrker kan tilegne sig, meget Andet, som han vil skære
bort som en Udvækst.

En saadan Udskillelse af det, man betragter som det Blivende,
fra det, som man anser for det Forgængelige hos Kierkegaard,
er Noget, som overhovedet slet ikke lader sig undgaa. Vil man
ingen saadan, da tage man Følgerne, og forkaste med Kierke\-%
gaard i hans sidste Optræden det hele kirkelige og det hele
menneskelige Samfundsliv, Ægteskabet osv., hvad dog vel Halv-%
% `Halv-/Grundtvigianerne': compound hyphen at a line end, kept (witness
%   2 reads `Halv-Grundtvigianerne').
Grundtvigianerne allermindst vil. Man pukke ikke paa, at man
er enig med ham i det Centrale, naar man selv for enhver Pris
vil afbryde den fineste og skarpeste Spids, i hvilken hans hele
Bestræbelse løber ud. Alt eller Intet, dette i vore Dage saa
yndede Raab, er et Løsen, der, naar Talen er om Anerkendelse,
ikke betegner Andet end et Attentat paa at tage Livet af al
Kritik.
\piecerule

% ============================================================
%  Fædrelandets Hr. 99 (printed pp. 102--105)
%  Indhold (PDF 9) reads: `Fædrelandets Hr. 99 102'
%  Head as printed: `«Fædrelandet»s Hr. 99 (Bjørnstjerne Bjørnson)' / `(1867)'
%    -- one italic line.  The only DOUBLE guillemets in the volume, and they
%    point OUTWARD like the single ones: « opens, » closes (confirmed at
%    600 dpi).  The closing parenthesis is italic.
%  Running head (pp. 103--105) reads: `Fædrelandets Hr. 99' (as Indhold).
%  One footnote, *) on p. 104, set small below a short rule; its source line
%    (`Sainte-Beuve: ...') stands on a line of its own, joined with \endgraf.
%  pp. 105--114 are printed lighter (ink density, not emphasis).
% ============================================================
% --- p. 102 ---
\piecehead{Fædrelandets Hr. 99}{\opage{102}\emph{«Fædrelandet»s Hr. 99 (Bjørnstjerne Bjørnson)}\\[2pt](1867)}
Den Ærede, der under Mærket 99 er optraadt imod mig i
\emph{Fædrelandet}, besidder blandt andre indtagende Egenskaber en
ikke almindelig Fripostighed. Det lyder utroligt, men det er
sandt, at Hr. 99 ikke engang selv har læst den Bog, i Anledning
af hvilken han er kommen i Harnisk. De Citater, han aftryk\-%
ker af denne Bog, findes aldeles ikke i Skriftet; men hvad han
fremsætter som Bogens Ord, det er \emph{Sætninger af mig}. Lige\-%
ledes: hvad han citerer af \emph{Fædrelandet}, er ikke Bladets Ord,
men \emph{mine}. Han anfører mig afvekslende med smaa Bogstaver
(som Bevis paa, hvor formastelig jeg er) og med store Bogstaver
(i den Formening, at han stiller mine Modstanderes Ord imod
mine); det forekommer mig, at dette lille Faktum i og for
sig er tilstrækkeligt til at betegne af hvad Art han som Modstan\-%
der er.

Ligesom sin Aandsfælle i \emph{Berlingske Tidende} springer 99
viselig alle Hovedpunkterne forbi og kaster sig over den ganske
uvæsenlige Jævnførelse, jeg har brugt. Jeg har sagt: Det er ikke
altid den stærkeste Krise, der hidfører det bedste og sandeste
Resultat; saaledes har f. Eks. i dette Aarhundredes Midte vi
Danske paa det politiske, Franskmændene paa det religions-%
filosofiske Omraade gjort et stort Skridt fremad uden stærke
% religions-/filosofiske: hyphen at the line end, KEPT: the volume prints
% both forms mid-line (hyphenated pp. 97, 98, 100; solid p. 108).  Witness 2
% reads `religions-filosofiske'.
Kriser, og omvendt har de politiske Revolutioner i Frankrig og
de teologisk-filosofiske Bevægelser i Danmark ikke umiddelbart
ført til Frihed. Dette var Jævnførelsen, og det forekommer mig,
% --- p. 103 ---
\opage{103}at den er klar. Man kan naturligvis hidse Aanden af den som
af enhver Sammenligning; men man er uberettiget til at ville
have den udstrakt videre, end jeg har strakt den; ti ellers er
det sandelig let nok at bringe den til at halte: ‹I den franske
Revolution udgødes Blod, i vore teologiske Stridigheder udgødes
intet, altsaa osv.›; ‹Napoleon den 3die har gjort Ende paa den
sidste Revolution, vor Videnskab har ingen Napoleon den 3die,
altsaa osv.› At vor politiske Bevægelse har den franske til
Forudsætning, medens Franskmændene Intet kender til Kierke\-%
gaard --- de kender meget vel den tyske Literatur --- er Jævn\-%
førelsen rent uvedkommende; men iøvrigt er denne Opfattelse af
Grunden til det Forhold, jeg har nævnt, saa flad som vel mu\-%
ligt. Grunden er den, at den romanske og den gotogermanske
Folkestamme har ganske forskellige Evner og Opgaver; den
romanske Stamme, repræsenteret af de Franske, opfatter og ud\-%
vikler Moralen socialt, den gotogermanske Stamme opfatter og
udvikler den personligt. Hin gør stadig Begyndelsen socialt, fra
denne udgaar Reformerne paa det personlige Omraade, og de
tilegner sig gensidigt Resultaterne af hinandens Bestræbelser. I
den Retning, i hvilken enhver af Stammerne er svagest udviklet,
faar den Resultaterne fra den anden uden megen Kamp, medens
den mere udviklede kan kæmpe sig saa træt, at den en Tid
lang ligger mat og ikke formaar at drage de Følger, den selv
har lagt tilrette. Hr. 99 stempler denne Anskuelse som ‹Non\-%
sens› og Jævnførelsen som en ‹usømmelig Fejltagelse›. Der kan
% Non-/sens: hyphen at the line end, joined as a break (`Nonsens');
% witness 2 reads `Non-sens'.  The volume prints `Nonsens' solid mid-line (p. 88 etc.), never hyphenated.
være saa stor en Kløft imellem to Menneskers Udvikling, at
den Enes Ord, om de er aldrig saa klare og vel begrundede,
nødvendigvis for den Anden maa tage sig ud som meningsløse.

Behøvedes der noget Bevis for min Ytring, at man her ikke
vil gendrive, men deklamere, saa har mine Modstandere nu
leveret det; ti hvormed har man besvaret mig uden med Inju\-%
rier, med Magtsprog og med Appel til den nationale Følelse.
Den Ene \emph{forbyder} mig at kalde Kierkegaard vor Filosofis Tycho
Brahe, Noget jeg imidlertid med eller uden hans Tilladelse
agter at fremture i. Den Anden vil overhovedet have mig al
Brug af Pen og Blæk forment, idetmindste i Dagbladene. O,
ædle Mand! Ja vist var det herligt, om man paa denne Maade
kunde skille sig af med de urolige Aander og sidde saa lunt i
sin egen Sæbekælder! Vist er det nemmere at kneble end at
% --- p. 104 ---
\opage{104}gendrive; men ak, hine Middelalderens gyldne Dage, da Sligt
gik saa glat og saa let --- de er forbi, de Festdage\fnmark{*)}{En
interessant Parallel med Hensyn til Slægtskabet mellem den
astronomiske og den filosofiske Forlæggelse af Centrum til Aabenbaringen
afgiver følgende Citat: ‹\emph{Pascal} har etsteds selv rettet en Sætning, taget
den tilbage og omformet den, for at faa Solen til at dreje sig om Jorden
og ikke Jorden om Solen. Denne store Aand, som paa dette Punkt sad
inde med en Rest af Overtro, veg tilbage for Kopernikus's Sandhed.›\endgraf
% The source line stands on its own indented line in the note; joined
% with \endgraf (a \par inside \fnmark breaks the compile).
Sainte-Beuve: Causeries du lundi III S. 254.}.

Men jeg kommer til det mest krænkende Angreb, der er
rettet paa mig. Det er det, at ‹den Mand, der kan føle og
tale› som jeg, ‹har intet Syn fælles med andre Danske, har
ingen dansk Bevidsthed›. Min Anskuelse er --- ikke ulogisk,
men --- \emph{udansk}. Hvad er det da, jeg har sagt? At Kierke\-%
gaard efter min Overbevisning har fejlet i det Centrale ɔ: i det,
som han antog for det Centrale, Læren om det Paradokse,
naturligvis ikke i det, jeg antager for det Centrale i hans
Skrifter, det rent humane Indhold. Er det en Forbrydelse
imod sit Fædreland at sige dette? Er hermed udtalt ‹Døds\-%
dommen over den danske Kultur›?› Har vi maaske faaet en
% PRINTER'S DEFECT: `Kultur›?›' -- a second closing › after the question
% mark with no opening mark to match it (600 dpi).  Transcribed as
% printed; the page's › count exceeds its ‹ count by one.
filosofisk Autoritet ovenikøbet ved Siden af de religiøse? Kierke\-%
gaard maa have Ret. Hvorfor? Fordi han er dansk. Og Mar\-%
tensen, hans Modstander? Ret, fordi han er dansk. Og Grundt\-%
vig, hans Antipode? Ret, fordi han er dansk. Og Tycho Brahe
imod Kopernikus? Ret, fordi han er dansk. Ja, det er bekendt
nok, at denne sidste Grundtvigske Følgeslutning allerede med
ophøjet Ringeagt for al sund Fornuft er bleven dragen. Ned
med den, som vover at sige, af den ‹danske Kultur› ikke i
% PRINTER'S DEFECT: `af den' for the sense `at den'; the page plainly
% prints `af' (600 dpi, beside `at sige').  Transcribed as printed.
alle Retninger er den bedste, sundeste, sandeste paa Jorden!

‹Vil \emph{Dagbladet} herefter holde Medarbejdere›, udraaber
Hr. 99 --- Udtrykket ‹holde› er allerede betegnende --- Med\-%
arbejdere, ‹som gentagende opmuntrer Danmarks Ungdom til
at gaa hid og se med egne Øjne, hvis den ikke vil kaldes
aandelig stavnsbunden›? Ja, er det ikke gyseligt at opfordre
den danske Ungdom til at se med egne Øjne? Nej, lad os
blive i det Spor, vi nu er, lad os tilhobe blive Nittengryn og
Nioghalvfemser; ti hvad hedder det ikke i Rævens \emph{Morale} hos
Holberg Art. 17 og Art. 18: ‹\emph{Scrupuleer} ikke udi Religionen;
% --- p. 105 ---
\opage{105}men holdt dig blindt ved den herskende Tro, saalænge som
den er paa Thronen, thi det er anstændigst og sikkerst at rette
sig efter Moden, saavel derudi som udi Klædedragt. --- Bemøye
dig ikke med at efterleede Sandhed; men troe hvad dig befales
at troe; thi det første er baade farligt og besværligt, det andet
derimod sikkert og mageligt; ja ofte baner Vey til \emph{Canonisation}›.
\piecerule

% ============================================================
%  R. Nielsens Lære om Tro og Viden (printed pp. 106--109)
%  Indhold (PDF 9) reads: `R. Nielsens Lære om Tro og Viden 106'
%  Head as printed: `Heegaard: Professor R. Nielsens Lære om' /
%    `Tro og Viden' / `(1867)' -- two italic lines, as set-up read them.
%  Running head (pp. 107--109) reads: `R. Nielsens Lære om Tro og Viden'
%    (as Indhold).
%  No footnotes.  Printed lighter throughout (ink density, not emphasis).
% ============================================================
% --- p. 106 ---
\piecehead{R. Nielsens Lære om Tro og Viden}{\opage{106}\emph{Heegaard: Professor R. Nielsens Lære om}\\
\emph{Tro og Viden}\\[2pt](1867)}
I disse Dage er udkommen \emph{Professor R. Nielsens Lære om
Tro og Viden}, en kritisk Afhandling af Dr. \emph{P. S. V. Heegaard}.
Den gør ingen Hemmelighed af, hvad den fører i sit Skjold;
ti allerede paa Bindet har den det Motto: ‹En Filosofi, der er
i aabenbar Modsigelse med den sunde Forstand, er falsk›.
Bogen er en literær Kuriositet, ikke saa meget ved sit Indhold,
ti en Del af hvad den forebringer, er sagt tilforn, som ved den
Overraskelse, det bereder, at se \emph{denne} Anskuelse udtalt fra \emph{denne}
Kant; en behagelig Overraskelse ganske vist, men derfor lige\-%
fuldt, som Nonpareil siger i \emph{Recensenten og Dyret}: la surprise
surprenante des surpris.

Forfatteren har nemlig som bekendt hørt til Professor
Nielsens Troeste, han har været dennes Lærling, Tilhænger og
Sekretær, og nu --- ak, hvor forandret! Hvor er nu den Tid,
da Dr. Heegaard som den højtbetroede Discipel laa ved Meste\-%
rens Bryst og kun rejste sig for med Strenghed at revse en
\emph{umoden} Ungdoms \emph{overilede} Forsøg paa at nedbryde en vis grund\-%
muret Lære? Ja, hvor er den Sne, som faldt ifjor? Og dog
havde Dr. Heegaard dengang Ret. Var det maaske ikke en
\emph{Overilelse?} Hvorfor skulde de unge Mennesker haste saa stærkt,
% `Overilelse?': the question mark is clearly italic (600 dpi), so inside.
hvorfor komme to hele Aar for tidligt? Kunde de ikke sagtens
vente lidt og tage Dr. Heegaard med; han vilde jo blot en lille
Tur udenlands først (se Bogens sidste Side). Og var det ikke
% --- p. 107 ---
\opage{107}et \emph{umodent} Foretagende at blæse til Angreb, inden Kærnetrop\-%
perne var indtrufne, at gaa paa uden at oppebi saa tapper og
værdifuld en Bistand? Jo visselig: ti hvorledes gik det? Ak, mise\-%
rabelt, Rækkerne vaklede, Angriberne var om en Hals; da, i det
afgørende Øjeblik, indtræffer Dr. Heegaard som Desaix ved
Marengo. Midt ud fra Professor Nielsens egen Lejr bryder han
paa én Gang med flyvende Faner og fuld Musik, efterladende
Intet uden en Basun, der efter sigende nu er ledig. Han haster
over mod den modsatte Side, han falder Professoren i Flanken;
det er sandelig Alvor, man bliver ganske højtideligt tilmode ved
saadant et Skue. Tillykke blander der sig en Smule Spøg og
Løjer i Optogets Højtidelighed og Krigens Gru.

For Eksempel: Det er den Paastand, siger Forfatteren, at
Professor Nielsen som en Aand af første Rang har iværksat et
stort filosofisk Gennembrud, hvilken har forarget ham og ægget
ham til at levere sin blodige Kritik. Dette er naturligvis en
Spøg; ti som bekendt er det Forfatteren selv, der tidligst og
ivrigst har sat denne Paastand i Stil. Fremdeles: det er højst
uheldigt, mener Dr. Heegaard, at Professor Nielsen, der begyndte
som ‹ren Hegelianer›, senere er bleven Hegels bestemte Mod\-%
stander. Atter Spøg, nemlig spøgefuld Illustration af Ordsproget:
Man skal ikke tale om Strikken osv. Endelig: Se til \emph{Grund\-%
ideernes Logik}, siger Forfatteren, det Bedste var, om Professor
Nielsen, der dog er ude af Stand til at bygge den færdig, straks
lod den staa hen ufuldendt som en ‹interessant Ruin›, altsaa som
en Slags Marmorkirke. ‹Selv Marmorkirke›, svarer Professor
Nielsen, rolig pegende paa første Del af den vidtberømte \emph{Indled\-%
ning til den rationelle Etik}. Er dette nu ikke en lystig Situation?

Dog Spøg tilside, det er jo klart nok, at netop den blan\-%
dede, halvt patetiske, halvt komiske Stilling, hvori Dr. Heegaard
som Forfatter af denne Bog befinder sig, i høj Grad forøger
Betydningen af hans kritiske Udtalelse. Med hvilken uimod\-%
staaelig Magt maa ikke Sandheden have vist sig for ham. siden
% PRINTER'S DEFECT: `ham. siden' -- a full stop where the sentence runs on
% (lower-case `siden'); plainly a stop at 600 dpi.  Transcribed as printed.
han har bragt den saa stort et Offer! Hans Kritik bærer da
ogsaa først og fremmest Vidnesbyrd om et sandhedskærligt Sind,
besluttet paa ikke mere at lade sig blænde, paa ikke at ville
lade sig nøje med Floskler eller Gaader, paa at ville blotte det
Væv af Spidsfindigheder og Fejlslutninger, paa hvilket Professor
Nielsens saakaldte Religionsfilosofi beror. Kritiken er ligesaa
dygtig som alvorlig og ligesaa skarp, som den er haardnakket
% --- p. 108 ---
\opage{108}og ubarmhjertig. Den sønderriver det skrøbelige Hjernespind, i
hvilket saa mange svage Sjæle er blevne fangne som Fluer, og
den er i sin Helhed saa afgørende og bestemt, at det nu
dog synes at maatte blive selv det langsomste Hoved klart,
at Professor Nielsen har tabt Slaget, og at hans Lære er gaaet
al Kødets Gang.

Det Nye ved Dr. Heegaards Behandlingsmaade af Spørgs\-%
maalet er især det, at han, istedenfor at vise, hvorledes Professor
Nielsens Lære er i Strid med Sandheden i og for sig, godtgør
dette ad en Omvej gennem en Paavisning af, hvorledes Nielsen
i alle Grene af sin Videnskab kommer i Modsigelse \emph{med sig selv}
ved sin fornuftstridige Opfattelse af Forholdet imellem Tro og
Viden. Til en Kritik af denne Natur var Dr. Heegaard med
sine særlige Forudsætninger, med sit udtømmende Kendskab til
Professor Nielsens hele System, netop fortrinlig rustet, og den
lader meget Lidet tilbage at ønske. Hist og her findes vel
nogen Uklarhed (som f. Eks. i Opfattelsen af Mysteriet som
Kategori) og nu og da en noget vrangvillig Fortolkning af
Modstanderens Ord (som hvor Mysteriet sættes lig med Menin\-%
gen), men disse Smaamangler forsvinder overfor den Grun\-%
dighed og Skarphed, som ellers udmærker Bevisførelsen Linje
for Linje.

Bogen er \emph{rent} kritisk, den angiver knap antydningsvis et
eget Standpunkt; Forfatteren synes netop lige at være kommen
ud af den Nielsenske Skole uden endnu selv at have vundet
Fodfæste. Det kan ikke negtes, at denne Omstændighed giver
Bogen og dens Stil noget Tørt og Forstandskoldt; der er i For\-%
fatterens Personlighed ingen Primitivitet, Intet af den Natur-%
Oprindelighed, der med Instinktets eller Genialitetens Sikkerhed
% `Natur-/Oprindelighed': compound hyphen at the line end (capital O),
% kept as printed.
griber det, der stiller sig som Sandhed for den; han synes som
Skribent at mangle det Bedste: den Frembringelsesgave, som
han efter ufrugtbare Naturers Vane vil fraskrive \emph{Tiden}.

Derfor undrer det os, at han ikke vurderer denne Evne
højere hos sin fordums Lærer. Ti Meget kan siges imod Pro\-%
fessor Nielsen, og Meget bør der siges imod ham: han har
viklet sig ind i et religionsfilosofisk Vilderede, og han er som
Tænker overhovedet ikke Systematiker; men han har den yppigt
frugtbare Aand, det geniale Blik, den sjældne Oprindelighed, der
er ligesaa væsenlige Evner hos en Tænker, som Tankens strengt
videnskabelige Dannelse og den skarpe og paalidelige Forstand.
% --- p. 109 ---
\opage{109}Man savner i Dr. Heegaards Bog den rette Anerkendelse heraf.
Den Taknemmelighed, som upaatvivlelig findes i Forfatterens
Sind imod hans saa højt begavede Lærer, kunde have fundet
et fyldigere Udtryk; den burde ikke have været indeklemt i en
kort Erklæring, men burde have gennemtrængt den hele Ud\-%
vikling og umuliggjort den opirrede Tone, i hvilken hans Kritik
adskillige Gange taler.
\piecerule

\end{document}
